\section*{Erd\H{o}s Problem 897: an explicit slowly growing counterexample}

\paragraph{Statement and result.}
An additive function means a function on positive integers
with \(f(ab)=f(a)+f(b)\) when \(\gcd(a,b)=1\).
We construct a positive such function, for \(n\geq2\), satisfying
\[
 \frac{f(p^k)}{\log(p^k)}\longrightarrow\infty
 \quad(p^k\longrightarrow\infty),
 \qquad
 \frac{f(n+1)-f(n)}{\log n}\longrightarrow0,
 \qquad
 \frac{f(n+1)}{f(n)}\longrightarrow1.
\]
Thus both proposed infinite-limsup conclusions are false.
The construction makes precise the slowly growing weights
credited to Erd\H{o}s and discussed by Wirsing; no novelty is claimed.

\paragraph{A slowly varying weight on prime powers.}
For real \(x>0\), let \(L(x)\) be the least nonnegative
integer \(j\) such that applying the natural logarithm \(j\)
times to \(x\) gives a number at most one. In particular,
\(L(x)=0\) for \(x\leq1\), and
\[
 L(x)=1+L(\log x)\quad(x>1).
\]
Set \(g(x)=\sqrt{L(x)+1}\). This is nondecreasing,
unbounded, and \(g(x)=O(\sqrt{\log x})\) as \(x\to\infty\).
For all sufficiently large \(x\), if \(t=\sqrt{\log x}\), then
\(t\geq\log\log x\), whence
\[
 L(t)\geq L(x)-2,\qquad
 0\leq g(x)-g(t)
 \leq\sqrt{L(x)+1}-\sqrt{L(x)-1}\longrightarrow0.
\tag{1}
\]
The jumps of \(L\) occur at the iterated-exponential thresholds
and these thresholds eventually have gaps greater than one.
Hence \(L(n+1)-L(n)\in\{0,1\}\) for all sufficiently large
integers \(n\), and \(g(n+1)-g(n)\to0\).

\paragraph{Definition and a uniform approximation.}
Put \(f(1)=0\), and, using the exact prime-power factors of \(n\), define
\[
 f(n)=\sum_{p^k\parallel n}g(p^k)\log(p^k).
\]
Unique prime factorization proves additivity on coprime arguments.
Also \(f(p^k)/\log(p^k)=g(p^k)\to\infty\).
We claim
\[
 f(n)=g(n)\log n+o(\log n).
\tag{2}
\]
Indeed \(\sum_{p^k\parallel n}\log(p^k)=\log n\).
Split this sum at \(t=\sqrt{\log n}\).
The factors exceeding \(t\) contribute at most
\((g(n)-g(t))\log n=o(\log n)\) to the nonnegative error
\(g(n)\log n-f(n)\), by (1).
There are at most \(\lfloor t\rfloor\) remaining factors,
each with logarithm at most \(\log t\). Their contribution
is at most \(g(n)t\log t=o(\log n)\).
For the last estimate one may use
\(L(n)=1+L(\log n)=O(\log\log n)\), which follows from
\(L(y)=O(\log y)\). Thus the error bound is uniform over
all factorizations of \(n\), proving (2).

\paragraph{Consecutive differences and ratios.}
The main term in (2) satisfies
\[
 \frac{g(n+1)\log(n+1)-g(n)\log n}{\log n}
 =g(n+1)-g(n)
  +g(n+1)\frac{\log(1+1/n)}{\log n}\longrightarrow0.
\]
Subtracting (2) at \(n\) and \(n+1\) gives the asserted
limit zero for the normalized difference.
Furthermore \(f(n)/\log n=g(n)+o(1)\to\infty\), so
\[
 \frac{f(n+1)}{f(n)}
 =1+\frac{f(n+1)-f(n)}{f(n)}\longrightarrow1.
\]
The ratio assertion concerns \(n\geq2\), where the denominator
is positive; \(f(1)=0\) causes no limiting difficulty.

\paragraph{Sources and scope.}
Current statement and historical attribution:
\url{https://www.erdosproblems.com/897}, checked 24 September 2026.
All estimates needed for this counterexample were proved above.
It addresses ordinary additive functions; it does not assert
the same result under the extra restrictions of strong or
complete additivity suggested in the historical discussion.
No local Lean verification is claimed.

\paragraph{Completion estimate.}
\textbf{COMPLETION: 100\%}
